% BHU PYQ -- Transcribed & Corrected
% Paper: BHUVA-12 / BHUVA12: Environmental Science (End Term)
% Exam: Under Graduate Semester I Examination, 2024-25 (NEP)
% Category: Value Added Course (VAC)
% Department: Environmental Science

\documentclass[11pt,a4paper]{article}
\usepackage[a4paper,top=2.5cm,bottom=2.5cm,left=2.8cm,right=2.8cm,headheight=14pt]{geometry}
\usepackage{amsmath,amssymb}
\usepackage[T1]{fontenc}
\usepackage[utf8]{inputenc}
\usepackage{lmodern,microtype,fancyhdr,enumitem,hyperref,lastpage,parskip}

\hypersetup{
    colorlinks=true,
    urlcolor=black,
    linkcolor=black,
    pdftitle={BHUVA12: Environmental Science -- BHU Sem I 2024-25 (NEP VAC)}
}

\pagestyle{fancy}
\fancyhf{}
\renewcommand{\headrulewidth}{0.4pt}
\renewcommand{\footrulewidth}{0.4pt}
\fancyhead[L]{\small   \textit{Banaras Hindu University}}
\fancyhead[R]{\small   \textit{BHUVA-12: Environmental Science (VAC)}}
\fancyfoot[C]{\small Page  \thepage\ of \pageref{LastPage}}


\newlist{parts}{enumerate}{2}
\setlist[parts,1]{label=(\Alph*),leftmargin=*,itemsep=3pt,topsep=2pt}

\newcommand{\pts}[1]{\hfill   \textit{[#1]}}


\begin{document}


\begin{center}
  {\large\bfseries BANARAS HINDU UNIVERSITY}\\[2pt]
  {
\normalsize Under Graduate --- Semester I Examination, 2024-25 (NEP)}\\[6pt]
  \rule{0.8\linewidth}{0.5pt}\\[6pt]
  {\large\bfseries Value Added Course (VAC)}\\[4pt]
  {\large\bfseries Paper Code: BHUVA-12 : Environmental Science (End Term)}\\[4pt]
  {
\normalsize Time: 1.5 Hours\hfill Maximum Marks: 70}\\[2pt]
  {\small REFNO. CE/Conf./D/231}
\end{center}

\rule{\linewidth}{0.4pt}

\smallskip



\noindent   \textit{Instructions: Answer all 70 Multiple Choice Questions. Each question carries equal marks (1 mark each).}

\bigskip


\noindent   \textbf{Question 1.} The major material used for making tools during the Neolithic Period was:\pts{1}

\begin{parts}
  \item Bronze
  \item Polished stone
  \item Steel
  \item Iron
\end{parts}

\medskip


\noindent   \textbf{Question 2.} The current human population across the globe is approximately:\pts{1}

\begin{parts}
  \item 6.9 billion
  \item 7 billion
  \item 9.5 billion
  \item 8.2 billion
\end{parts}

\medskip


\noindent   \textbf{Question 3.} Anthropocentrism states that:\pts{1}

\begin{parts}
  \item All living beings have intrinsic value
  \item Human beings are the central or most significant entities in the world
  \item Ecosystems have legal rights
  \item Non-human entities have superior rights
\end{parts}

\medskip


\noindent   \textbf{Question 4.} The main objective of environmental education is to:\pts{1}

\begin{parts}
  \item Increase industrial growth
  \item Promote environmental awareness and sustainable action
  \item Promote urbanization
  \item Reduce agricultural production
\end{parts}

\medskip


\noindent   \textbf{Question 5.} Biomagnification refers to:\pts{1}

\begin{parts}
  \item Increase in concentration of toxic substances up the food chain
  \item Increase in population of a species
  \item Growth of plants using fertilizers
  \item Increase in atmospheric temperature
\end{parts}

\medskip


\noindent   \textbf{Question 6.} In India, the agency which is responsible for monitoring and controlling pollution at the national level is:\pts{1}

\begin{parts}
  \item NITI Aayog
  \item Central Pollution Control Board (CPCB)
  \item Supreme Court
  \item Ministry of Earth Sciences
\end{parts}

\medskip


\noindent   \textbf{Question 7.} The lowermost layer of the atmosphere where we live is called:\pts{1}

\begin{parts}
  \item Mesosphere
  \item Troposphere
  \item Stratosphere
  \item Thermosphere
\end{parts}

\medskip


\noindent   \textbf{Question 8.} The term "Biosphere" refers to:\pts{1}

\begin{parts}
  \item The zone of Earth where life exists
  \item The gaseous envelope surrounding Earth
  \item The solid rock layer of Earth
  \item All oceans and freshwater bodies
\end{parts}

\medskip


\noindent   \textbf{Question 9.} Which of the following is a primary air pollutant?\pts{1}

\begin{parts}
  \item Ozone
  \item Carbon monoxide
  \item Smog
  \item Peroxyacetyl nitrate (PAN)
\end{parts}

\medskip


\noindent   \textbf{Question 10.} Deforestation leads to:\pts{1}

\begin{parts}
  \item Increase in soil fertility
  \item Loss of biodiversity and soil erosion
  \item Decrease in atmospheric carbon dioxide
  \item Increase in rainfall
\end{parts}

\medskip


\noindent   \textbf{Question 11.} An ecosystem consists of:\pts{1}

\begin{parts}
  \item Only biotic components
  \item Only abiotic components
  \item Interactions between biotic and abiotic components
  \item Only decomposers
\end{parts}

\medskip


\noindent   \textbf{Question 12.} What do you mean by the term pollution?\pts{1}

\begin{parts}
  \item The natural process of ecosystem balance
  \item The undesirable change in the environment due to harmful substances
  \item The increase in biodiversity in an area
  \item The sustainable use of natural resources
\end{parts}

\medskip


\noindent   \textbf{Question 13.} Acid rain is primarily caused by emissions of:\pts{1}

\begin{parts}
  \item Carbon dioxide and methane
  \item Sulfur dioxide and nitrogen oxides
  \item Chlorofluorocarbons
  \item Oxygen and nitrogen
\end{parts}

\medskip


\noindent   \textbf{Question 14.} Ozone depletion can lead to:\pts{1}

\begin{parts}
  \item Increased risk of skin cancer due to UV radiation
  \item Global cooling
  \item Increased soil nitrogen content
  \item Reduction in ocean salinity
\end{parts}

\medskip


\noindent   \textbf{Question 15.} COVID-19 is an example of:\pts{1}

\begin{parts}
  \item Pandemic
  \item Epidemic
  \item Endemic
  \item None of the above
\end{parts}

\medskip


\noindent   \textbf{Question 16.} Which of the following is a major cause of soil pollution?\pts{1}

\begin{parts}
  \item Overuse of chemical fertilizers
  \item Application of cow dung
  \item Drip irrigation
  \item Planting more trees
\end{parts}

\medskip


\noindent   \textbf{Question 17.} Which of the following is an example of a natural disaster?\pts{1}

\begin{parts}
  \item Earthquake
  \item Oil spill
  \item Nuclear accident
  \item Gas leak
\end{parts}

\medskip


\noindent   \textbf{Question 18.} The concept of Sustainable Development emphasizes:\pts{1}

\begin{parts}
  \item Meeting present needs without compromising future generations
  \item Maximizing resource exploitation today
  \item Stopping all industrial activities
  \item Focusing only on economic growth
\end{parts}

\medskip


\noindent   \textbf{Question 19.} Greenhouse effect is primarily caused by:\pts{1}

\begin{parts}
  \item Trapping of heat by greenhouse gases in the atmosphere
  \item Depletion of topsoil
  \item Excessive rainfall
  \item Cooling of ocean currents
\end{parts}

\medskip


\noindent   \textbf{Question 20.} Which of the following is a non-renewable energy source?\pts{1}

\begin{parts}
  \item Solar energy
  \item Coal
  \item Wind energy
  \item Hydroelectric power
\end{parts}

\medskip


\noindent   \textbf{Question 21.} The term "Ecosystem" was coined by:\pts{1}

\begin{parts}
  \item Arthur Tansley
  \item Charles Darwin
  \item Ernst Haeckel
  \item Eugene Odum
\end{parts}

\medskip


\noindent   \textbf{Question 22.} What is the main reason for air pollution?\pts{1}

\begin{parts}
  \item Industrial emissions and vehicular exhaust
  \item Tree plantation
  \item Use of electric vehicles
  \item Use of microbial fertilizers
\end{parts}

\medskip


\noindent   \textbf{Question 23.} Which of the following is a major threat to biodiversity?\pts{1}

\begin{parts}
  \item Over-exploitation of natural resources and habitat loss
  \item Tree plantation drives
  \item Organic farming
  \item Wildlife sanctuary creation
\end{parts}

\medskip


\noindent   \textbf{Question 24.} Ramsar Convention is an international treaty for the conservation of:\pts{1}

\begin{parts}
  \item Wetlands
  \item Deserts
  \item Forests
  \item Coral reefs
\end{parts}

\medskip


\noindent   \textbf{Question 25.} Chipko Movement in India was aimed at:\pts{1}

\begin{parts}
  \item Protecting trees from being cut down
  \item Cleaning rivers
  \item Reducing plastic usage
  \item Wildlife protection
\end{parts}

\medskip


\noindent   \textbf{Question 26.} Eutrophication in water bodies is caused by excess of:\pts{1}

\begin{parts}
  \item Nutrients like nitrogen and phosphorus
  \item Heavy metals like lead
  \item Dissolved oxygen
  \item Salt content
\end{parts}

\medskip


\noindent   \textbf{Question 27.} What is the main reason for water pollution?\pts{1}

\begin{parts}
  \item Interlinking of rivers
  \item Sustainable agricultural practices
  \item Industrial effluents and untreated sewage
  \item Floods
\end{parts}

\medskip


\noindent   \textbf{Question 28.} Which of the following is a decomposer in an ecosystem?\pts{1}

\begin{parts}
  \item Plants
  \item Elephant
  \item Crow
  \item Microorganisms (Bacteria and Fungi)
\end{parts}

\medskip


\noindent   \textbf{Question 29.} Which of the following is a renewable source of energy?\pts{1}

\begin{parts}
  \item Coal
  \item Natural gas
  \item Solar energy
  \item Petroleum
\end{parts}

\medskip


\noindent   \textbf{Question 30.} Which layer of the atmosphere contains the ozone layer?\pts{1}

\begin{parts}
  \item Troposphere
  \item Stratosphere
  \item Mesosphere
  \item Thermosphere
\end{parts}

\medskip


\noindent   \textbf{Question 31.} Minamata disease was caused by environmental poisoning of:\pts{1}

\begin{parts}
  \item Mercury
  \item Lead
  \item Cadmium
  \item Arsenic
\end{parts}

\medskip


\noindent   \textbf{Question 32.} Among the following, which one constitutes the major proportion of Earth's interior by volume?\pts{1}

\begin{parts}
  \item Crust
  \item Mantle
  \item Core
  \item Lithosphere
\end{parts}

\medskip


\noindent   \textbf{Question 33.} The possible consequence of unplanned urbanization is:\pts{1}

\begin{parts}
  \item Generation of excessive waste
  \item Increase in air pollution
  \item Social inequality
  \item All of the above
\end{parts}

\medskip


\noindent   \textbf{Question 34.} Biodiversity hotspot refers to a region with:\pts{1}

\begin{parts}
  \item High level of endemic species under threat
  \item Low species richness
  \item Only marine life
  \item Extinct species
\end{parts}

\medskip


\noindent   \textbf{Question 35.} Which gas is predominantly responsible for global warming?\pts{1}

\begin{parts}
  \item Carbon dioxide (CO2)
  \item Oxygen (O2)
  \item Nitrogen (N2)
  \item Argon (Ar)
\end{parts}

\medskip


\noindent   \textbf{Question 36.} Noise pollution is measured in units of:\pts{1}

\begin{parts}
  \item Hertz (Hz)
  \item Decibel (dB)
  \item Watt (W)
  \item Joule (J)
\end{parts}

\medskip


\noindent   \textbf{Question 37.} The United Nations Sustainable Development Goals (SDGs) consist of:\pts{1}

\begin{parts}
  \item 15 goals
  \item 16 goals
  \item 17 goals
  \item 18 goals
\end{parts}

\medskip


\noindent   \textbf{Question 38.} Which of the following is an essential role of soil?\pts{1}

\begin{parts}
  \item Provides nutrients and water for plants
  \item Powers winds and ocean currents
  \item Filters air for human survival
  \item Provides renewable energy source
\end{parts}

\medskip


\noindent   \textbf{Question 39.} What is the main characteristic of a non-renewable resource?\pts{1}

\begin{parts}
  \item It replenishes naturally at a very fast rate
  \item It gets exhausted upon continuous usage
  \item It can never be depleted
  \item It produces zero waste
\end{parts}

\medskip


\noindent   \textbf{Question 40.} Which instrument is used to measure earthquake intensity/magnitude?\pts{1}

\begin{parts}
  \item Barometer
  \item Seismograph / Richter Scale
  \item Thermometer
  \item Anemometer
\end{parts}

\medskip


\noindent   \textbf{Question 41.} Which of the following is an example of habitat fragmentation?\pts{1}

\begin{parts}
  \item Creation of new wildlife reserves
  \item Construction of roads and dams through forests
  \item Restoration of wetlands
  \item Afforestation programs
\end{parts}

\medskip


\noindent   \textbf{Question 42.} Which type of ecosystem is characterized by permafrost and treeless landscapes?\pts{1}

\begin{parts}
  \item Grasslands
  \item Deserts
  \item Tundra
  \item Wetlands
\end{parts}

\medskip


\noindent   \textbf{Question 43.} What is the primary consequence of thermal pollution in water bodies?\pts{1}

\begin{parts}
  \item Decrease in dissolved oxygen levels
  \item Increase in fish growth rate
  \item Purification of water
  \item Freezing of water surface
\end{parts}

\medskip


\noindent   \textbf{Question 44.} Which environmental movement in Karnataka was inspired by the Chipko Movement?\pts{1}

\begin{parts}
  \item Appiko Movement
  \item Narmada Bachao Andolan
  \item Silent Valley Movement
  \item Jungle Bachao Andolan
\end{parts}

\medskip


\noindent   \textbf{Question 45.} Kyoto Protocol is an international agreement focused on:\pts{1}

\begin{parts}
  \item Reducing greenhouse gas emissions
  \item Banning hazardous chemical trade
  \item Protecting endangered birds
  \item Disposal of electronic waste
\end{parts}

\medskip


\noindent   \textbf{Question 46.} Solid waste management strategy "3 Rs" stands for:\pts{1}

\begin{parts}
  \item Reduce, Reuse, Recycle
  \item Read, Research, Report
  \item Remove, Replace, Restore
  \item Refine, React, Rebuild
\end{parts}

\medskip


\noindent   \textbf{Question 47.} What is the primary role of the National Green Tribunal (NGT) in India?\pts{1}

\begin{parts}
  \item Promote primary education
  \item Handle environmental disputes and enforce legal rights
  \item Encourage rapid industrialization
  \item Reduce poverty
\end{parts}

\medskip


\noindent   \textbf{Question 48.} The United Nations Environment Programme (UNEP) headquarters is located in:\pts{1}

\begin{parts}
  \item Geneva
  \item Nairobi
  \item New York
  \item Paris
\end{parts}

\medskip


\noindent   \textbf{Question 49.} Primary consumers in a food chain are also known as:\pts{1}

\begin{parts}
  \item Herbivores
  \item Carnivores
  \item Omnivores
  \item Decomposers
\end{parts}

\medskip


\noindent   \textbf{Question 50.} Carrying capacity of an environment refers to:\pts{1}

\begin{parts}
  \item Maximum population size that an environment can sustain indefinitely
  \item Total weight of soil in an area
  \item Speed of ocean currents
  \item Maximum rainfall recorded
\end{parts}

\medskip


\noindent   \textbf{Question 51.} Earth Day is celebrated globally every year on:\pts{1}

\begin{parts}
  \item 15th May
  \item 22nd May
  \item 22nd April
  \item 1st April
\end{parts}

\medskip


\noindent   \textbf{Question 52.} The national animal of India is:\pts{1}

\begin{parts}
  \item The Royal Bengal Tiger
  \item Lion
  \item Elephant
  \item One-horned Rhinoceros
\end{parts}

\medskip


\noindent   \textbf{Question 53.} The LiFE (Lifestyle for Environment) Mission launched by India is about:\pts{1}

\begin{parts}
  \item Providing insurance to economically weaker sections
  \item Adopting nature-friendly lifestyle for environmental protection
  \item Literacy mission
  \item Farmer loan waiver
\end{parts}

\medskip


\noindent   \textbf{Question 54.} The lithosphere is:\pts{1}

\begin{parts}
  \item The solid, outermost rocky layer of the Earth
  \item The water component of the Earth
  \item The gaseous envelope of the Earth
  \item The oceanic trench part
\end{parts}

\medskip


\noindent   \textbf{Question 55.} Which of the following is an urban development mission in India?\pts{1}

\begin{parts}
  \item AMRUT Mission
  \item HRIDAY Mission
  \item Smart Cities Mission
  \item All of the above
\end{parts}

\medskip


\noindent   \textbf{Question 56.} Which of the following is a major reason for human migration?\pts{1}

\begin{parts}
  \item Better employment opportunities and living conditions
  \item Absence of technology
  \item High pollution levels
  \item Voluntary isolation
\end{parts}

\medskip


\noindent   \textbf{Question 57.} World Environment Day is observed on:\pts{1}

\begin{parts}
  \item 5th June
  \item 11th July
  \item 1st December
  \item 16th September
\end{parts}

\medskip


\noindent   \textbf{Question 58.} World Ozone Day is celebrated on:\pts{1}

\begin{parts}
  \item 16th September
  \item 22nd March
  \item 5th June
  \item 22nd April
\end{parts}

\medskip


\noindent   \textbf{Question 59.} Which of the following is a water-borne disease?\pts{1}

\begin{parts}
  \item Cholera
  \item Arthritis
  \item Common cold
  \item Cough
\end{parts}

\medskip


\noindent   \textbf{Question 60.} Which of the following is a clear indicator of global climate change?\pts{1}

\begin{parts}
  \item Melting of glaciers and ice-caps
  \item Decreasing ocean levels
  \item Monotone weather patterns
  \item Decreasing global average temperature
\end{parts}

\medskip


\noindent   \textbf{Question 61.} The direct and indirect benefits that humans derive from ecosystems are called:\pts{1}

\begin{parts}
  \item Ecosystem services
  \item Ecosystem structure
  \item Ecosystem regeneration
  \item Ecosystem degradation
\end{parts}

\medskip


\noindent   \textbf{Question 62.} Photovoltaic cells convert solar energy directly into:\pts{1}

\begin{parts}
  \item Electrical energy
  \item Chemical energy
  \item Thermal energy
  \item Nuclear energy
\end{parts}

\medskip


\noindent   \textbf{Question 63.} In an ecological pyramid, producers belong to which trophic level?\pts{1}

\begin{parts}
  \item First Trophic Level
  \item Second Trophic Level
  \item Third Trophic Level
  \item Fourth Trophic Level
\end{parts}

\medskip


\noindent   \textbf{Question 64.} Which of the following is an example of ex-situ conservation?\pts{1}

\begin{parts}
  \item Botanical Gardens and Seed Banks
  \item National Parks
  \item Biosphere Reserves
  \item Wildlife Sanctuaries
\end{parts}

\medskip


\noindent   \textbf{Question 65.} Project Tiger was launched in India in the year:\pts{1}

\begin{parts}
  \item 1973
  \item 1986
  \item 1992
  \item 2001
\end{parts}

\medskip


\noindent   \textbf{Question 66.} In a food web, the maximum energy resides at the level of:\pts{1}

\begin{parts}
  \item Primary consumers
  \item Secondary consumers
  \item Tertiary consumers
  \item Producers
\end{parts}

\medskip


\noindent   \textbf{Question 67.} Which international organization publishes the Red List of Threatened Species?\pts{1}

\begin{parts}
  \item IUCN
  \item Ministry of Environment, Forest & Climate Change
  \item Central Pollution Control Board
  \item Ministry of Earth Sciences
\end{parts}

\medskip


\noindent   \textbf{Question 68.} Wild Life Protection Act in India was enacted in:\pts{1}

\begin{parts}
  \item 1972
  \item 1980
  \item 1986
  \item 1998
\end{parts}

\medskip


\noindent   \textbf{Question 69.} Which gas released during the Bhopal Gas Tragedy in 1984 caused thousands of deaths?\pts{1}

\begin{parts}
  \item Methyl Isocyanate (MIC)
  \item Phosgene
  \item Carbon Monoxide
  \item Mustard Gas
\end{parts}

\medskip


\noindent   \textbf{Question 70.} The main objective of Project Elephant launched in 1992 in India is to:\pts{1}

\begin{parts}
  \item Protect wild elephants and their natural habitats/corridors
  \item Domesticate elephants for transport
  \item Export elephants to zoos
  \item Promote elephant safaris
\end{parts}

\medskip

\vfill
\rule{\linewidth}{0.4pt}

\begin{center}\small
\href{https://bhu-exam.vercel.app/}{Click here to visit our website}\\[4pt]
\href{https://me-aryan.vercel.app/}{Click here to visit our contributor}
\end{center}

\end{document}
